\chapter{Cài đặt hệ thống}

\section{Môi trường phát triển}
\begin{table}[H]
\centering
\begin{tabularx}{\textwidth}{|X|X|}
\hline
\textbf{Thành phần} & \textbf{Phiên bản} \\
\hline
Hệ điều hành & Windows \\
\hline
Trình biên dịch & MSVC (Microsoft Visual Studio Compiler) / GCC \\
\hline
Hệ thống Build & CMake (version 3.20 trở lên) \\
\hline
Thư viện xử lý ảnh & OpenCV 4.x \\
\hline
Ngôn ngữ lập trình & C++17 \\
\hline
\end{tabularx}
\caption{Môi trường và công cụ phát triển dự án.}
\end{table}

\section{Tiền xử lý dữ liệu}
Trong chương trình, các ảnh đầu vào đều trải qua một số bước tiền xử lý cố định thông qua module \texttt{retrieval\_utils}:
\begin{itemize}
    \item \textbf{Đọc ảnh:} Ảnh được đọc trực tiếp qua hàm \texttt{cv::imread}. Nếu file hỏng, bỏ qua.
    \item \textbf{Đồng bộ kích thước:} Hầu hết ảnh được resize về $256 \times 256$ để giảm nhiễu cục bộ và tăng tốc tìm kiếm. Riêng thuật toán HOG được resize cố định về $64 \times 64$ để phù hợp với chuẩn rút trích.
    \item \textbf{Khử nhiễu:} Áp dụng Gaussian Blur với Kernel kích thước $5 \times 5$.
    \item \textbf{Chuyển không gian màu:} Tùy phương pháp, ảnh sẽ được đổi qua RGB (cho SIFT, ORB), HSV (cho Color Histogram) hoặc Grayscale (cho LBP, HOG).
\end{itemize}

\section{Quy trình truy vấn Top-K thực tế}
Dưới đây là luồng xử lý do \texttt{ImageSearchEngine::search} đảm nhận, tự động định tuyến tìm kiếm qua các chỉ mục tương ứng:
\begin{minted}[
    fontsize=\small,
    breaklines=true,
    linenos=true,
    frame=single,
    tabsize=4
]{cpp}
// 1. Nhận ảnh đầu vào và rút trích đặc trưng
cv::Mat queryFeat = extractor->extract(queryImage);

// === Nhánh 1: Sử dụng Inverted Index cho các thuật toán BoW (SIFT/ORB) ===
if (m_invertedIndices.find(method) != m_invertedIndices.end() &&
    m_invertedIndices.at(method).isBuilt()) {
    
    auto& invIdx = m_invertedIndices.at(method);
    auto invResults = invIdx.search(queryFeat, db.size());
    
    // Chuyển kết quả TF-IDF sang SearchResult (đảo chiều score)
    float maxScore = invResults.empty() ? 0.f : invResults.front().second;
    for (const auto& [imgID, score] : invResults) {
        results.push_back({db[imgID].imagePath, db[imgID].label, maxScore - score});
    }
}
// === Nhánh 2: Sử dụng FLANN KD-Tree cho Euclidean/Cosine (HOG, Color_HOG) ===
else if ((metric == DistanceMetric::EUCLIDEAN || metric == DistanceMetric::COSINE) &&
         m_flannIndices.find(method) != m_flannIndices.end()) {
         
    auto flannIndex = m_flannIndices.at(method);
    cv::Mat queryF32;
    queryFeat.convertTo(queryF32, CV_32F);
    
    flannIndex->knnSearch(queryF32, indices, dists, db.size(), cv::flann::SearchParams(32));
    for (int i = 0; i < db.size(); i++) {
        int idx = indices.at<int>(0, i);
        float dist = dists.at<float>(0, i);
        float finalScore = (metric == DistanceMetric::EUCLIDEAN) ? std::sqrt(dist) : dist / 2.0f;
        results.push_back({db[idx].imagePath, db[idx].label, finalScore});
    }
}
// === Nhánh 3: Sử dụng Forward Index cho Chi-Squared (Color Hist, LBP) ===
else {
    if (m_forwardIndex.isBuilt()) {
        for (int imgID = 0; imgID < db.size(); ++imgID) {
            const cv::Mat& dbFeat = m_forwardIndex.getFeature(imgID, method);
            float score = computeDistance(queryFeat, dbFeat, metric);
            const auto& meta = m_forwardIndex.getMetadata(imgID);
            results.push_back({meta.imagePath, meta.label, score});
        }
    }
}

// Sắp xếp tăng dần theo score (score nhỏ = giống nhất)
std::sort(results.begin(), results.end());
\end{minted}

\section{Cải tiến tốc độ và Hiệu năng}
Chương trình triển khai các giải pháp cải thiện tốc độ trực tiếp:
\begin{itemize}
    \item \textbf{Lưu trữ dữ liệu Offline:} Quá trình trích xuất đặc trưng cực kỳ tốn chi phí. Thay vì trích xuất toàn bộ thư mục mỗi lần chạy lại phần mềm, tiến trình \texttt{DatabaseBuilder} sẽ rút trích trước 100\% tập dữ liệu và kết xuất ra các file \texttt{.yml}.
    \item \textbf{Tối ưu tra cứu bằng Forward Index:} Lưu trữ đặc trưng trực tiếp bằng hash map giúp truy xuất đặc trưng của bất kỳ ảnh nào trong $O(1)$ thay vì phải tìm kiếm tuần tự trong các map lồng nhau của từng bản ghi dữ liệu phẳng.
    \item \textbf{Tìm kiếm thưa qua Inverted Index:} Đối với các đặc trưng BoW thưa, Inverted Index giúp bỏ qua việc so khớp với các ảnh không liên quan, tăng tốc độ tính toán gấp nhiều lần. Sử dụng mảng vector tuyến tính để tích lũy điểm TF-IDF thay vì map động giúp tăng tối đa tốc độ truy cập bộ nhớ đệm (cache-friendly).
    \item \textbf{Chỉ mục FLANN (Fast Library for Approximate Nearest Neighbors):} Để tăng tốc thuật toán K-Nearest Neighbors với các vector dense (HOG), cấu trúc cây phân nhánh KD-Tree được xây dựng. Thuật toán tìm kiếm xấp xỉ này cho tốc độ tra cứu gần như tức thời.
\end{itemize}

\section{Xử lý ngoại lệ}
\begin{table}[H]
\centering
\begin{tabularx}{\textwidth}{|l|X|}
\hline
\textbf{Trường hợp} & \textbf{Cách xử lý} \\
\hline
Không chọn ảnh & Chương trình thông báo yêu cầu chọn ảnh lên màn hình giao diện. \\
\hline
Ảnh bị lỗi & \texttt{imread} báo rỗng, thuật toán in ra cảnh báo trên Terminal và tự động bỏ qua. \\
\hline
Thiếu file \texttt{.yml} & \texttt{GUIApp} khởi động tính năng cảnh báo và dừng lại yêu cầu chạy \texttt{build\_features.exe} trước. \\
\hline
Tập CSDL bé hơn $K$ & Hàm trả về \texttt{min(K, database.size())}. \\
\hline
\end{tabularx}
\caption{Xử lý các ngoại lệ hệ thống.}
\end{table}
